\section{Truncation and Effective Dimension}
\label{sec:ch46}

Sections~\ref{sec:ch43}--\ref{sec:ch45} identified radiative rank \(r_{\mathrm{rad}}(\varepsilon)\), spectral compression, and
geometry-capped coupling on \(\boldsymbol{\Gamma}_\omega\) \cite{stratton1941,harrington2001,devaney2004}. Section~\ref{sec:ch46}
replaces informal cutoffs by explicit angular-order ceilings, scaling exponents, coupling-tail laws, and the effective
electromagnetic dimension \(d_{\mathrm{eff}}\) that laboratories report as ``number of modes'' \cite{hansen1988}. Finite support
fixes the exponents; representation choice does not \cite{jackson1999,stratton1941}.

\subsection{Maximum supported angular order}
\label{sec:ch461}
% TOC tag: [HOME][USE SS3.7.2]

Boundary quantization in Subsection~\ref{sec:ch372} replaces continuous direction spectra by countable mode sets under
\(\Gamma_{b}\) \cite{stratton1941,harrington2001}. On open regions with compact sources, the analogous realizability question is:
which \(\ell\) sectors carry export gain above tolerance on \(\mathcal{J}_{\mathrm{real}}\)?
\cite{hansen1988,devaney2004}.

\paragraph{Assumptions.}
Fix ball support \(B_{a}\), flux-normalized VSH listing \(\mathcal{C}\), and tolerance \(\varepsilon\in(0,1)\)
\cite{stratton1941,hansen1988}. Presuppose Eq.~\eqref{eq:ch4.high-ell-decay} and Definition~\ref{def:ch4.radiative-rank}
\cite{harrington2001,devaney2004}.

\paragraph{Maximum supported order.}
Define the \emph{maximum supported angular order}
\begin{equation}
\label{eq:ch4.ell-max-supported}
\ell_{\mathrm{supp}}(\varepsilon,k_{0}a)
=
\max\left\{
\ell:\ \exists\,(\ell,m,\sigma)\ \text{with}\ 
\sigma_{(\ell,m,\sigma)}\ge\varepsilon\,\sigma_{1}
\ \text{in Eq.~\eqref{eq:ch4.gamma-SVD}}
\right\},
\end{equation}
with \(\sigma_{(\ell,m,\sigma)}\) the singular value aligned to \(\widehat{\mathbf{u}}_{\ell m}^{(\sigma)}\) in the export basis
\cite{harrington2001,hansen1988}. Equation~\eqref{eq:ch4.ell-max-supported} is the synthesis-side angular cutoff: modes with
\(\ell>\ell_{\mathrm{supp}}\) lie below export threshold at \(\varepsilon\) regardless of listing labels
\cite{stratton1941,devaney2004}.

\begin{theorem}[Supported order versus electrical size]
\label{thm:ch4.ell-max-scaling}
For compact \(B_{a}\) and aligned VSH gauge,
\begin{equation}
\label{eq:ch4.ell-max-scaling}
\ell_{\mathrm{supp}}(\varepsilon,k_{0}a)
\ \le\
\ell_{\max}(\eta_{\mathrm{tail}},k_{0}a)
\ \sim\
k_{0}a+\mathcal{O}\!\big(\log(1/\varepsilon)\big),
\end{equation}
importing Eq.~\eqref{eq:ch4.ellmax-from-tail} and Eq.~\eqref{eq:ch4.high-ell-decay} \cite{hansen1988,jackson1999,stratton1941}.
Equality holds in the isotropic ball case up to gauge constants \cite{harrington2001,devaney2004}.
\end{theorem}
\begin{proof}
Tail envelope Eq.~\eqref{eq:ch4.high-ell-decay} bounds nonzero \(\widehat{c}_{\ell m}^{(\sigma)}\); singular-value ordering of
\(\boldsymbol{\Gamma}_\omega\) on VSH sectors gives the cutoff identity \cite{hansen1988,stratton1941,harrington2001}.
\end{proof}

Theorem~\ref{thm:ch4.ell-max-scaling} is the open-region counterpart of Subsection~\ref{sec:ch372} quantization: geometry
imposes a countable ceiling on \(\ell\), not merely on mesh indices \cite{stratton1941}. Phase-only metasurfaces on fixed \(a\)
cannot raise \(\ell_{\mathrm{supp}}\) without raising \(k_{0}a\) or feed rank \cite{devaney2004,hansen1988}.

\begin{figure}[t]
\centering
\ChOneFigBlock{%
\begin{tikzpicture}[ch1fig, node distance=7mm and 10mm]
  \node[ch1box] (ka) {$k_{0}a$};
  \node[ch1box, right=12mm of ka] (ell) {$\ell_{\mathrm{supp}}$};
  \node[ch1box, right=12mm of ell] (sr) {$\mathcal{S}_{\mathrm{real}}$};
  \draw[ch1flow] (ka) -- (ell) -- (sr);
\end{tikzpicture}%
}
\caption{Subsection~\ref{sec:ch461}: electrical size sets maximum supported \(\ell\) before synthesis listing.}
\label{fig:ch4.ell-max-supported}
\end{figure}

\paragraph{Worked illustration (handset aperture).}
At \(k_{0}a=2\), \(\ell_{\mathrm{supp}}(10^{-2})\approx4\)--\(6\) by Eq.~\eqref{eq:ch4.ell-max-scaling}; marketing claims of
dominant \(\ell=12\) OAM on the same aperture fail Theorem~\ref{thm:ch4.ell-max-scaling} \cite{hansen1988,jackson1999,harrington2001}.

\paragraph{Problem (supported order audit).}
\textbf{Problem.} A VSH listing retains \(\ell\le15\) but \(k_{0}a=1.8\). How many \(\ell\) sectors matter at \(\varepsilon=10^{-2}\)?

\textbf{Step 1.} Compute \(\ell_{\mathrm{supp}}(10^{-2},1.8)\) from Eq.~\eqref{eq:ch4.ell-max-supported} via SVD of \(\mathbf{T}\)
\cite{harrington2001}.

\textbf{Step 2.} Compare to Eq.~\eqref{eq:ch4.ell-max-scaling} for conservative bound \cite{hansen1988,jackson1999}.

\textbf{Step 3.} Truncate design targets to \(\ell\le\ell_{\mathrm{supp}}\) \cite{devaney2004,stratton1941}.

\textbf{Physical meaning.} \(\ell_{\mathrm{supp}}\) is a geometry-imposed ceiling, not a plotting convenience \cite{hansen1988}.

\paragraph{Validity.}
Eqs.~\eqref{eq:ch4.ell-max-supported}--\eqref{eq:ch4.ell-max-scaling} assume open-region ball supports and aligned VSH gauge
\cite{stratton1941,hansen1988}. Cavities require Subsection~\ref{sec:ch372} indices \((\ell,m,n_{r})\) \cite{harrington2001}.


\subsection{Scaling laws for spectral truncation}
\label{sec:ch462}
% TOC tag: [HOME]

Truncation exponents summarize how fast \(\mathcal{S}_{\mathrm{real}}(\mathcal{C})\) grows with electrical size and frequency
\cite{stratton1941,jackson1999,hansen1988}. Subsection~\ref{sec:ch462} states the scaling laws as realizability inequalities on
\(r_{\mathrm{rad}}(\varepsilon)\) and \(\ell_{\mathrm{supp}}\) \cite{harrington2001,devaney2004}.

\paragraph{Assumptions.}
Presuppose Eqs.~\eqref{eq:ch4.weyl-mode-count}, \eqref{eq:ch4.geometry-coupling-scaling}, and
Theorem~\ref{thm:ch4.support-volume-scaling} \cite{stratton1941,hansen1988}. Asymptotic regime \(k_{0}a\ge1\) unless noted
\cite{jackson1999}.

\paragraph{Volume and area exponents.}
On \(B_{a}\),
\begin{equation}
\label{eq:ch4.truncation-volume-scaling}
r_{\mathrm{rad}}(\varepsilon)
\ \sim\
C_{\mathrm{vol}}(\varepsilon)\,(k_{0}a)^{3},
\qquad
\ell_{\mathrm{supp}}(\varepsilon,k_{0}a)
\ \sim\
C_{\ell}(\varepsilon)\,k_{0}a,
\end{equation}
importing Eq.~\eqref{eq:ch4.weyl-mode-count} and Theorem~\ref{thm:ch4.ell-max-scaling} \cite{stratton1941,jackson1999,hansen1988}.
On aperture \(\Sigma_{a}\) of area \(A_{a}\),
\begin{equation}
\label{eq:ch4.truncation-area-scaling}
\operatorname{rank}\mathbf{K}\big|_{\Sigma_{a}}
\ \sim\
C_{A}(\varepsilon)\,k_{0}^{2}A_{a},
\end{equation}
from Eq.~\eqref{eq:ch4.aperture-coupling-scaling} \cite{collin2001,harrington2001,devaney2004}. Volume scaling dominates rank growth
when the enclosing ball is the synthesis domain; area scaling governs aperture-only models \cite{stratton1941}.

\begin{proposition}[Scaling law consistency]
\label{prop:ch4.truncation-scaling-consistency}
If equivalence Theorem~\ref{thm:ch4.volume-surface-equivalence} holds and \(\mathcal{D}_{\mathrm{cpl}}=1\),
Eqs.~\eqref{eq:ch4.truncation-volume-scaling} and \eqref{eq:ch4.truncation-area-scaling} bound the same \(\mathbf{c}\) image up to
constants depending on \(\varepsilon\) and mesh alignment \cite{stratton1941,harrington2001,collin2001}. Neither exponent is raised by
phase-only retuning on fixed geometry \cite{devaney2004,hansen1988}.
\end{proposition}

\paragraph{Worked illustration (doubling frequency).}
At fixed \(a\), doubling \(k_{0}\) doubles \(k_{0}a\) in Eq.~\eqref{eq:ch4.truncation-volume-scaling}: \(r_{\mathrm{rad}}\) grows
roughly by \(2^{3}=8\) in the large-body regime \cite{jackson1999,harrington2001,hansen1988}.

\paragraph{Problem (volume versus area rank).}
\textbf{Problem.} A patch array on \(\Sigma_{a}\) reports rank \(k_{0}^{2}A_{a}\approx50\); enclosing ball gives \((k_{0}a)^{3}\approx30\).
Which bound governs synthesis?

\textbf{Step 1.} Identify synthesis domain: aperture-only versus enclosed volume \cite{collin2001,stratton1941}.

\textbf{Step 2.} Apply the bound matching the declared \(\mathcal{J}_{\mathrm{real}}\) support \cite{harrington2001,devaney2004}.

\textbf{Step 3.} Use the smaller conservative rank for realizability audits \cite{hansen1988}.

\textbf{Physical meaning.} Scaling exponents attach to the declared support class \cite{stratton1941}.

\paragraph{Validity.}
Eqs.~\eqref{eq:ch4.truncation-volume-scaling}--\eqref{eq:ch4.truncation-area-scaling} are asymptotic; Rayleigh regime \(k_{0}a\ll1\)
requires dipole-dominated rank \cite{jackson1999,hansen1988}.


\subsection{Decay of high-order coupling}
\label{sec:ch463}
% TOC tag: [HOME][USE SS3.5.3]

Mode density and coupling tails in Subsection~\ref{sec:ch353} count angular channels on \(S^{2}\) \cite{jackson1999,hansen1988}.
Subsection~\ref{sec:ch463} restates the tail as a realizability law on \(\boldsymbol{\mathcal{K}}_{\mathrm{geo}}\) from
Eq.~\eqref{eq:ch4.geometry-coupling-operator} \cite{stratton1941,harrington2001}.

\paragraph{Assumptions.}
Fix compact \(\Omega_{s}\), flux gauge, and coupling operator \(\boldsymbol{\mathcal{K}}_{\mathrm{geo}}\)
\cite{stratton1941,hansen1988}. Import Eq.~\eqref{eq:ch4.coupling-ell-decay} \cite{harrington2001,devaney2004}.

\paragraph{Coupling tail envelope.}
For \(\ell>k_{0}D_{s}\),
\begin{equation}
\label{eq:ch4.coupling-tail-envelope}
\sup_{\Jimp\in\mathcal{J}_{\mathrm{real}}}
\frac{|\mathcal{K}_{\ell m}^{(\sigma)}[\Jimp]|}{\|\boldsymbol{\mathcal{K}}_{\mathrm{geo}}[\Jimp]\|}
\le
\mathcal{C}_{\mathrm{tail}}(\ell,k_{0}D_{s})\,
\ell^{-1}\exp(-\ell/k_{0}D_{s}),
\end{equation}
extending Eq.~\eqref{eq:ch4.coupling-ell-decay} with explicit sector constants \(\mathcal{C}_{\mathrm{tail}}\)
\cite{hansen1988,stratton1941,jackson1999}. Equation~\eqref{eq:ch4.coupling-tail-envelope} ties Subsection~\ref{sec:ch353} channel
counts to export gain: high-\(\ell\) channels may be listable yet carry negligible coupling \cite{harrington2001,devaney2004}.

\begin{theorem}[Coupling tail implies spectral compression]
\label{thm:ch4.coupling-tail-compression}
If Eq.~\eqref{eq:ch4.coupling-tail-envelope} holds, then
\(r_{\mathrm{rad}}(\varepsilon)\le C(k_{0}D_{s},\varepsilon)\) with \(C\) growing slower than channel count
\((2\ell_{\max}+1)\) from Subsection~\ref{sec:ch353} \cite{hansen1988,stratton1941,harrington2001}. Listing size therefore
overstates synthesis dimension when \(\ell_{\max}\gg k_{0}D_{s}\) \cite{devaney2004}.
\end{theorem}
\begin{proof}
Summing envelope Eq.~\eqref{eq:ch4.coupling-tail-envelope} over \((m,\sigma)\) and comparing to \(\varepsilon\sigma_{1}\) in
Definition~\ref{def:ch4.radiative-rank} \cite{hansen1988,harrington2001,stratton1941}.
\end{proof}

\paragraph{Worked illustration (high channel count, low rank).}
At \(\ell_{\max}=20\) and \(k_{0}D_{s}=1.5\), Subsection~\ref{sec:ch353} lists hundreds of labels but
\(r_{\mathrm{rad}}(10^{-2})\approx10\) by Theorem~\ref{thm:ch4.coupling-tail-compression} \cite{hansen1988,harrington2001,devaney2004}.

\paragraph{Problem (coupling versus DOS).}
\textbf{Problem.} Channel DOS grows as \(\ell_{\max}^{2}\) but measured diversity saturates. Reconcile.

\textbf{Step 1.} Distinguish listable channels from export couplings above \(\varepsilon\sigma_{1}\) \cite{hansen1988,stratton1941}.

\textbf{Step 2.} Apply Theorem~\ref{thm:ch4.coupling-tail-compression} \cite{harrington2001,devaney2004}.

\textbf{Step 3.} Report \(r_{\mathrm{rad}}(\varepsilon)\), not raw DOS \cite{jackson1999}.

\textbf{Physical meaning.} DOS counts labels; coupling tails count watts \cite{hansen1988}.

\paragraph{Validity.}
Eq.~\eqref{eq:ch4.coupling-tail-envelope} is narrowband and assumes isotropic exterior \cite{stratton1941}. Anisotropic feeds modify
\(\mathcal{C}_{\mathrm{tail}}\) without removing decay \cite{harrington2001}.


\subsection{Effective rank and electromagnetic dimension}
\label{sec:ch464}
% TOC tag: [HOME]

Radiative rank names export directions; laboratories require a single integer \(d_{\mathrm{eff}}\) for arrays, MIMO, and inverse design
\cite{harrington2001,devaney2004}. Subsection~\ref{sec:ch464} locks that integer to singular-value thresholds on
\(\boldsymbol{\Gamma}_\omega\) \cite{stratton1941,hansen1988}.

\paragraph{Assumptions.}
Presuppose Definition~\ref{def:ch4.radiative-rank} and Eq.~\eqref{eq:ch4.effective-dimension-preview}
\cite{harrington2001,devaney2004}. Tolerance \(\varepsilon\) declared for synthesis and certification jointly \cite{hansen1988}.

\begin{definition}[Effective electromagnetic dimension]
\label{def:ch4.effective-em-dimension}
The \emph{effective electromagnetic dimension} of \((\Omega_{s},\omega,\mathcal{C})\) at tolerance \(\varepsilon\) is
\begin{equation}
\label{eq:ch4.effective-em-dimension}
d_{\mathrm{eff}}(\varepsilon)
=
r_{\mathrm{rad}}(\varepsilon)
=
\#\big\{n:\ \sigma_{n}\ge\varepsilon\,\sigma_{1}\big\},
\end{equation}
with \(\{\sigma_{n}\}\) from Eq.~\eqref{eq:ch4.gamma-SVD} \cite{stratton1941,harrington2001,hansen1988}.
\end{definition}

Equation~\eqref{eq:ch4.effective-em-dimension} elevates Eq.~\eqref{eq:ch4.effective-dimension-preview} to the Chapter~\ref{ch:04} name for
geometry-imposed synthesis dimension \cite{devaney2004}. \(d_{\mathrm{eff}}\) is not patch count, not \(\ell_{\max}\), and not channel DOS
\cite{hansen1988,harrington2001}.

\begin{theorem}[Dimension bounds realizability]
\label{thm:ch4.dimension-bounds-realizability}
\(\dim\mathcal{S}_{\mathrm{real}}(\mathcal{C})\le d_{\mathrm{eff}}(\varepsilon)\) at tolerance \(\varepsilon\), and every
\(\mathbf{c}\in\mathcal{S}_{\mathrm{real}}(\mathcal{C})\) has at most \(d_{\mathrm{eff}}(\varepsilon)\) independent components above
export threshold \cite{stratton1941,harrington2001,devaney2004,hansen1988}.
\end{theorem}
\begin{proof}
Theorem~\ref{thm:ch4.radiative-rank} and Definition~\ref{def:ch4.effective-em-dimension} \cite{harrington2001,stratton1941}.
\end{proof}

\begin{figure}[t]
\centering
\ChOneFigBlock{%
\begin{tikzpicture}[ch1fig, node distance=7mm and 9mm]
  \node[ch1box] (svd) {Eq.~\eqref{eq:ch4.gamma-SVD}};
  \node[ch1box, right=11mm of svd] (rr) {$r_{\mathrm{rad}}$};
  \node[ch1box, right=11mm of rr] (deff) {$d_{\mathrm{eff}}$};
  \node[ch1box, right=11mm of deff] (sr) {$\mathcal{S}_{\mathrm{real}}$};
  \draw[ch1flow] (svd) -- (rr) -- (deff) -- (sr);
\end{tikzpicture}%
}
\caption{Subsection~\ref{sec:ch464}: effective electromagnetic dimension is radiative rank in the export gauge.}
\label{fig:ch4.effective-em-dimension}
\end{figure}

\paragraph{Worked illustration (MIMO claim).}
A \(16\times16\) patch array claims \(256\) synthesis DOF; \(d_{\mathrm{eff}}(10^{-2})=22\) on the enclosing ball at \(k_{0}a=3\)
\cite{harrington2001,collin2001,hansen1988}. MIMO rank must be reported as \(d_{\mathrm{eff}}\), not element count \cite{devaney2004}.

\paragraph{Problem (dimension from MoM).}
\textbf{Problem.} SVD of \(\mathbf{T}\) shows 180 singular values above \(10^{-3}\). Is \(d_{\mathrm{eff}}(10^{-3})=180\)?

\textbf{Step 1.} Compare to continuum \(\boldsymbol{\Gamma}_\omega\) per Lemma~\ref{lem:ch4.rank-gap-aliasing} \cite{harrington2001}.

\textbf{Step 2.} Use \(d_{\mathrm{eff}}=\min(\text{mesh rank},r_{\mathrm{rad}})\) in aligned gauge \cite{devaney2004,hansen1988}.

\textbf{Step 3.} Report \(\varepsilon\) and \(k_{0}D_{s}\) with \(d_{\mathrm{eff}}\) \cite{stratton1941}.

\textbf{Physical meaning.} \(d_{\mathrm{eff}}\) is the geometry-imposed synthesis dimension \cite{hansen1988}.

\paragraph{Validity.}
Definition~\ref{def:ch4.effective-em-dimension} depends on \(\mathcal{C}\) normalization; \(\varepsilon\) must be stated with \(\sigma_{1}\)
\cite{hansen1988,stratton1941}. Loss lowers tail decay and may increase \(d_{\mathrm{eff}}\) at fixed geometry \cite{harrington2001}.

Section~\ref{sec:ch46} has replaced heuristic truncation by \(\ell_{\mathrm{supp}}\), scaling exponents, coupling-tail compression, and
\(d_{\mathrm{eff}}\) \cite{stratton1941,hansen1988,harrington2001}. Section~\ref{sec:ch47} treats null spaces, accessibility, and
reconstruction limits on the same operator chain \cite{devaney2004}.


\section{Null Spaces, Accessibility, and Reconstruction}
\label{sec:ch47}

Export rank \(d_{\mathrm{eff}}\) counts synthesizable directions; null spaces and finite observation rank determine what can be
reconstructed, certified, or inverted from measured data \cite{stratton1941,harrington2001,devaney2004}. Section~\ref{sec:ch47}
develops that split: non-radiating and unreachable modes, reactive accessibility, inversion limits, and localization versus angular
resolution \cite{hansen1988}.
